\documentclass[11pt]{article}
\usepackage[T1]{fontenc}
\usepackage{lmodern,amsmath,amssymb,amsthm,mathtools,enumitem,hyperref}
\usepackage[margin=1in]{geometry}
\newtheorem{theorem}{Theorem}[section]
\newtheorem{lemma}[theorem]{Lemma}
\newtheorem{proposition}[theorem]{Proposition}
\newtheorem{corollary}[theorem]{Corollary}
\theoremstyle{definition}\newtheorem{definition}[theorem]{Definition}
\newcommand{\norm}[1]{\nu(#1)}
\newcommand{\subc}[2]{#1\langle #2\rangle}
\newcommand{\lit}{\mathrel{\equiv}}
\newcommand{\Rfree}{\mathcal R_{\mathrm{free}}}
\newcommand{\Cnorm}{\mathcal C_\nu}
\newcommand{\Knorm}{\mathcal K_\nu}
\newcommand{\ipra}{\mathsf{iPRA}}
\newcommand{\file}[1]{{\small\nolinkurl{#1}}}
\title{Canonical successor-seeded syntax:\\a six-rule presentation, with its normalization cost}
\author{Working mathematical review}
\date{November 11, 2026}
\begin{document}
\maketitle

\begin{abstract}
A successor-seeded scalar recursor can be presented with only its successor computation as a proof rule if its zero computation is imposed as a canonical-syntax convention. Combining this convention with the established crossed-transport and witness-directed-introduction constructions gives six inference patterns on terms with no independent successor constructor. We prove the new normalization and derivation equivalences, including the required body-extensionality argument. The arithmetic bridge is the previously proved successor-seeded bridge, whose full bootstrap and native finite-case repair are reproduced in the appendices. The removed computation clause must be executed by normalized substitution and charged there. The accompanying complete report supplies a matching strict checker comparison that includes this cost; the mathematical proof in this source establishes the change of presentation.
\end{abstract}

\section{Scope and the exact mathematical change}

The comparison source is ProveIt commit
\texttt{8eb71c1e570e07ef45ead686628e77f18ef34798}.
The successor-seeded recursor and its arithmetic completeness are inherited from
\file{Research/successor_seed_completion/} and
\file{Research/scalar_successor/}.
The crossed-transport argument is the retraction criterion of
\file{Research/reversal_core/theory/reversal_equivalence.tex}.
The single-introduction argument is the idea formalized in
\file{Lean/ConciseIPRA/NativeOneIntro.lean}.
The new construction here is their application to explicitly canonical successor-seeded syntax, with a proof of the resulting change of term equality and substitution.

No independent successor constructor is used. Raw terms are
\[
t::=x\mid 0\mid R(n,a;u.h).
\]
Only \(u\) in \(h\) is bound. Put
\[
Sa:=R(0,a;u.0),\qquad 1:=S0.
\]
Bound variables are renamed to avoid capture. There are no numerical variables over function descriptions: a body is finite earlier syntax, with its remaining free variables as passive parameters.

\section{Canonical terms and substitution}

Define a normalization of raw terms recursively:
\[
\norm{x}=x,\qquad \norm{0}=0,
\]
\[
\norm{R(n,a;u.h)}=
\begin{cases}
R(0,\norm a;u.0),&n\lit0,\\
R(\norm n,\norm a;u.\norm h),&n\not\lit0.
\end{cases}
\]
Here \(n\lit0\) means literal syntactic identity with zero. It is not a numerical equality test. A canonical term is a fixed point of \(\nu\). Equivalently, every subterm with literal zero counter has literal zero body. This is a decidable grammar of finite expressions.

The normalizer preserves the outer term constructor:
\[
\norm t\lit0\quad\Longleftrightarrow\quad t\lit0.
\]
In particular \(\norm{Sa}=S\norm a\). It is idempotent, respects renaming of bound and free variables, and introduces no new free variables.

Extend \(\nu\) to native formulas and contexts by normalizing all terms:
\[
A::=(s=t)\mid Q_y A\mid D_{x,y}A.
\]
The roots \(=\), \(Q\), and \(D\) are preserved. All terms in a canonical formula are canonical.

\begin{lemma}[Substitution coherence]\label{cr:coherence}
For every raw term or formula \(T\) and capture-free simultaneous numerical substitution \(\sigma\),
\[
\norm{T[\sigma]}=
\norm{\,\norm T[\nu\circ\sigma]\,}.
\]
\end{lemma}
\begin{proof}
Use structural induction after freshening bound names. The only additional case is an \(R\)-node. If its original counter is literally zero, both sides erase the body. If it is not literally zero but becomes zero under substitution, both sides again erase the body. In the remaining case both sides retain the \(R\)-node and the equality follows from the induction hypotheses for counter, initializer, and body. The fact that normalization never turns a nonzero-root constructor into literal zero justifies the case division. Formula binders introduce no extra case.
\end{proof}

For canonical expressions define
\[
\subc T\sigma:=\norm{T[\sigma]}.
\]
This is the substitution operation used by every rule below. Lemma~\ref{cr:coherence} gives associativity of normalized substitution and its compatibility with renaming. A variable substituted by zero may expose a zero-counter node. Normalizing only the original syntax, or testing the unsubstituted counter, is insufficient.

In a canonical-calculus display, an expression \(R(n,a;u.h)\) denotes the normalization of that raw constructor expression. This smart-constructor convention is relevant when \(n\lit0\): its written body is discarded. In the raw comparison calculus below, constructor expressions instead retain their literal bodies. Substitution may accept a raw image because it normalizes the completed result; equivalently it may first normalize that image by Lemma~\ref{cr:coherence}.

The sole primitive computation schema is
\[
R(Sn,a;u.h)=\subc h{\norm{R(n,a;u.h)}/u}.
\tag{\(R_+\)}\label{cr:rplus}
\]
All displayed terms are canonical: the left counter \(Sn\) is never literally zero. The raw zero computation
\(R(0,a;u.h)=Sa\) is now definitional, because both sides have the same canonical spelling.

\section{The complete six-rule calculus}

Contexts are finite reusable lists of canonical native formulas. Define root reversal by
\[
(s=t)^\dagger=(t=s),\qquad
(Q_yB)^\dagger=Q_yB,\qquad
(D_{x,y}B)^\dagger=D_{x,y}B.
\]
It commutes with canonical substitution.

The calculus \(\Knorm\) has the following six inference patterns.

\begin{enumerate}[leftmargin=*]
\item \textbf{Hypothesis:} \(\Gamma\vdash A\) when \(A\in\Gamma\).
\item \textbf{Successor computation:} every instance of \eqref{cr:rplus}, in every context.
\item \textbf{Crossed equational transport:}
\[
\frac{\Gamma\vdash s=t\qquad
      \Gamma\vdash \subc E{s/v}}
     {\Gamma\vdash (\subc E{t/v})^\dagger},
\qquad s\not\lit0.
\tag{X}
\]
Here \(E(v)\) is a literal-equation template. Constant templates are allowed.
\item \textbf{Crossed closure elimination:}
\[
\frac{\Gamma\vdash Q_yB(y)\qquad\Gamma\vdash\subc B{0/y}}
     {\Gamma\vdash(\subc B{t/y})^\dagger}.
\tag{QX}
\]
\item \textbf{Paired elimination:}
\[
\frac{\Gamma\vdash D_{x,y}B(x,y)\qquad
      \Gamma,Q_y\subc B{b/x}\vdash A}
     {\Gamma\vdash A}.
\tag{DE}
\]
The eigenparameter \(b\) is fresh for the retained context, conclusion, and fixed parameters of \(B\).
\item \textbf{One introduction:}
\[
\frac{\Gamma,\subc B{t/x,\eta/y}
          \vdash\subc B{t/x,Sz/y}}
     {\Gamma\vdash D_{x,y}B},
\qquad
\eta=
\begin{cases}
z,&B\text{ is a literal equation and }t\lit0,\\
0,&\text{otherwise}.
\end{cases}
\tag{WI}
\]
Here \(z\) is fresh for the retained context, witness \(t\), and fixed parameters of \(B\).
\end{enumerate}

There is no primitive reflexivity, symmetry, numerical case rule, induction rule, zero-computation rule, ordinary logical connective, or general numerical-substitution rule. The designated contradiction is \(0=S0\).

The intended ordinary readings are
\[
(Q_yB)^\circ=B^\circ(0)\to\forall y\,B^\circ(y),
\]
\[
(D_{x,y}B)^\circ=
\exists x\bigl(B^\circ(x,0)\to\forall y\,B^\circ(x,y)\bigr).
\]
All numerical terms use the successor-seeded primitive-recursive interpretation of \(R\). Domains are whole, and equality is actual equality.

\section{First establish the unrestricted comparison}

Do not initially attempt a substitution theorem for the six-rule calculus. Literal-zero source and witness tests are not stable under arbitrary numerical substitution. Instead use two unrestricted comparison calculi.

\begin{definition}[Raw and canonical unrestricted calculi]
The raw calculus \(\Rfree\) uses raw \(R\)-terms, ordinary capture-free substitution, the two computation schemas
\[
R(0,a;u.h)=Sa,\qquad
R(Sn,a;u.h)=h[R(n,a;u.h)/u],
\]
Hyp, unrestricted transport in every native formula template, ordinary QE, DE, and the two introductions
\[
\frac{\Gamma,B(t,0)\vdash B(t,Sz)}
     {\Gamma\vdash D_{x,y}B},
\qquad
\frac{\Gamma,E(t,z)\vdash E(t,Sz)}
     {\Gamma\vdash D_{x,y}E}
\quad(E\text{ a literal equation}).
\]
The canonical calculus \(\Cnorm\) has the same unrestricted rules using canonical substitution, but has only \(R_+\) as its computation schema.

Here ordinary QE is QX without root reversal. Unrestricted transport is X without either its literal-zero source restriction or its root reversal, and it permits arbitrary native formula templates. The freshness conditions for DE and both introductions are the conditions already stated for DE and WI. Hypothesis and DE are unchanged.
\end{definition}

\begin{lemma}[Structural rules]\label{cr:struct}
Weakening, eigenvariable renaming, ordinary-assumption cut, and capture-free numerical substitution are admissible in \(\Rfree\) and \(\Cnorm\).
Weakening, eigenvariable renaming, and ordinary-assumption cut are admissible in \(\Knorm\), without presupposing numerical substitution there.
\end{lemma}
\begin{proof}
Induct on derivations, freshening local eigenparameters first. For \(\Cnorm\), coherence of normalized substitution proves closure of the computation schema and all displayed instances. The case/progress distinction in the unrestricted calculi depends only on the formula root and is stable. Cut replaces assumption uses with proofs without changing witnesses, roots, or source terms; hence it remains valid in the restricted six-rule calculus. Renaming preserves literal zero and root constructors. This proof makes no assertion that general numerical substitution directly preserves an X or WI node.
\end{proof}

\begin{lemma}[Elementary resources]\label{cr:resources}
Both unrestricted calculi derive reflexivity, symmetry, transitivity, numerical congruence, ordinary \(Q\)-introduction by cases, numerical cases, and contextual induction on literal equations.
\end{lemma}
\begin{proof}
Define the retraction context
\[
J(a):=R(S0,0;u.a).
\]
The successor computation gives \(J(a)=a\), even in \(\Cnorm\). Use this equation twice in transport with motive \(v=a\) to get \(a=a\). The usual equality motives give symmetry, transitivity, and numerical congruence.

From \(\Gamma,B(0)\vdash B(Sz)\), case introduction with an unused witness gives \(D_{w,y}B(y)\); DE followed by Hyp exports \(Q_yB(y)\). Applying QE with supplied base and successor proofs gives numerical cases. Repeating the unused-witness construction with progress introduction gives equation induction from its ordinary base and step premises. No selector, predecessor, or ordinary logical connective is needed.
\end{proof}

\begin{lemma}[Recursor body extensionality]\label{cr:body}
Suppose, for a fresh numerical variable \(v\), that a fixed context proves \(h(v)=k(v)\). Then it proves
\[
R(n,a;u.h(u))=R(n,a;u.k(u)).
\]
The assertion holds in both unrestricted calculi.
\end{lemma}
\begin{proof}
Apply contextual literal-equation induction to the equation comparing the two trajectories at a fresh counter. At zero both terms are \(Sa\), by zero computation in \(\Rfree\), and by canonical syntax and reflexivity in \(\Cnorm\). At the successor, computation exposes the two updates. Use the induction equation to replace the old state in \(h\), then use the supplied pointwise body equality at the second old state. Such instantiation is justified by Lemma~\ref{cr:struct}, which has already been proved for these unrestricted calculi. All passive parameters and the initializer remain fixed during the induction.
\end{proof}

This lemma is necessary. Changing a normalized body changes its finite recursor description; it is not generally ordinary congruence of one fixed first-order function symbol.

\begin{lemma}[Normalization is provable in the raw calculus]\label{cr:normal-eq}
For every raw term \(t\),
\[
\Rfree\vdash t=\norm t.
\]
For every raw native formula \(A\), the raw calculus derives both
\(A\vdash\norm A\) and \(\norm A\vdash A\), uniformly in its free parameters.
\end{lemma}
\begin{proof}
For terms use structural induction. Counter and initializer changes use numerical congruence. The body change uses Lemma~\ref{cr:body} and the induction hypothesis with the body's bound variable fresh and free in the auxiliary proof. If the normalized counter is zero, use the raw zero computation to replace the body by zero; otherwise the resulting term is already its normal form.

For formulas use external induction on the native formula syntax. The equation case follows from the two endpoint identities. For the \(Q\) case, under \(Q_yA(y)\), introduce \(Q_y\norm{A(y)}\) by cases: convert its assumed normalized base to the old base, use QE at the fresh successor, and convert back. The reverse direction is identical. For \(D\), eliminate the old \(D\) at a fresh witness, introduce the new \(D\) at that witness by cases, convert the new base to the old base, use the exported old \(Q\), and convert the successor conclusion. Cut closes the construction. These are explicit derived binder-congruence arguments, not applications of an assumed ordinary quantifier logic.
\end{proof}

\begin{theorem}[Canonical and raw consequence]\label{cr:canonical-equiv}
For canonical \(\Gamma,A\),
\[
\Gamma\vdash_{\Cnorm}A
\quad\Longleftrightarrow\quad
\Gamma\vdash_{\Rfree}A.
\]
More generally, normalization sends every raw derivation to a canonical derivation of its normalized sequent.
\end{theorem}
\begin{proof}
For the raw-to-canonical direction normalize the entire finite derivation. Coherence of substitution handles every instance of transport, QE, DE, and either introduction. A raw zero-computation leaf normalizes to \(Sa=Sa\), which is derived by Lemma~\ref{cr:resources}; a successor-computation leaf normalizes to an \(R_+\) instance.

For the converse, regard canonical expressions as raw expressions and recurse on the canonical derivation. Any canonical rule instance differs from its raw counterpart only by normalizations of substituted expressions. Lemma~\ref{cr:normal-eq} converts these premises, local assumptions, and conclusions in both directions. Use ordinary-assumption cut for local assumptions. For example, canonical transport first converts the canonical minor \(\nu(E[s/v])\) to raw \(E[s/v]\), applies raw transport, and converts \(E[t/v]\) to its canonical normalization. Introductions and DE use the same conversions under their retained local assumptions. Freshness is preserved because normalization introduces no free variables. No restricted-rule substitution theorem enters either direction.
\end{proof}

\section{Only then reduce the canonical calculus to six rules}

\subsection{Equality before introduction repair}

In \(\Knorm\), keep \(J(a)=R(S0,0;u.a)\) and \(c=J(0)\). Each \(J(a)\) is literally nonzero and \(J(a)=a\) is an \(R_+\) instance.

X with major \(c=0\) and a constant equation template reverses any proved equation. Compound root reversal is the identity. Therefore QX followed by reversal supplies ordinary QE, and X followed by reversal supplies ordinary equational transport at nonzero literal sources.

Reflexivity follows by using \(J(a)=a\) as both X premises with template \(v=a\). To remove the literal-zero source restriction, use the established retraction argument, whose short details are included here.

If \(p,q\not\lit0\), \(p=0\), and \(q=0\), then X with major \(p=0\), minor \(J(p)=p\), and template \(J(v)=p\) yields \(p=c\). Similarly \(q=c\). Reverse and join through the allowed source \(c\) to obtain \(p=q\). Transitivity now holds without restriction: if its middle term is nonzero use restricted ordinary transport; if it is zero and both endpoints are nonzero, reverse the second equation and use the zero join; literal-zero endpoint cases are immediate.

Numerical congruence along a nonzero source is ordinary restricted transport from reflexivity. For an equation \(0=t\) with \(t\not\lit0\), reverse it to \(t=0\); transport the reflexivity \(f(t)=f(t)\) in the template \(f(v)=f(t)\), giving \(f(0)=f(t)\). If both endpoints are zero, use reflexivity. Finally, from \(s=t\) and \(a(s)=b(s)\), numerical congruence and unrestricted transitivity give
\[
a(t)=a(s)=b(s)=b(t).
\]
Thus unrestricted ordinary transport in literal-equation templates is available before any introduction repair. All steps instantiate explicit schemas at the terms at hand; none substitutes into a restricted derivation.

\subsection{Recovering both introductions}

\begin{lemma}[Case introduction]\label{cr:case}
Every unrestricted case-introduction instance is derivable in \(\Knorm\).
\end{lemma}
\begin{proof}
WI already is case introduction when its matrix is compound or its witness is not literally zero. Only an equation matrix \(E\) at witness zero remains.

Use instead the witness \(c=J(0)\), which is nonliteral and provably zero. Under the local assumption \(E(c,0)\), unrestricted equation transport gives \(E(0,0)\). Cut in the given case-premise derivation to get \(E(0,Sz)\); transport back to \(E(c,Sz)\). WI at witness \(c\) is the case choice and concludes the required unchanged formula \(D_{x,y}E(x,y)\). The new witness contains no eigenparameter \(z\).
\end{proof}

\begin{lemma}[Progress introduction]\label{cr:progress}
Every unrestricted literal-equation progress instance is derivable in \(\Knorm\).
\end{lemma}
\begin{proof}
Suppose the supplied premise is
\(\Gamma,E(t,z)\vdash E(t,Sz)\).
Let the first coordinate \(w\) be unused in the equation \(E(t,y)\). WI with witness zero proves
\[
\Gamma\vdash D_{w,y}E(t,y)
\]
by its progress branch. Under the additional assumption \(E(t,0)\), eliminate this paired formula. Its exported clause is \(Q_yE(t,y)\), independent of the fresh witness. Ordinary QE gives \(E(t,Sz)\). Closing DE gives the premise required for the recovered case introduction at the original witness \(t\), which concludes \(D_{x,y}E(x,y)\).
\end{proof}

\subsection{Compound transport}

Now equation transport, ordinary QE, DE, and both introductions are available. Derive ordinary transport in all native formula templates by external induction on their constructor structure. The equality case is already proved. For a \(Q\)-template, use derived case \(Q\)-introduction: convert its new base back along the equality by the induction hypothesis, apply the old \(Q\) at the fresh successor, and convert forward. For a \(D\)-template, DE chooses the old witness; case-introduce the new \(D\) at that witness, convert its base back, use the exported old clause, and convert the successor forward. Reverse the numerical equality to obtain both directions when required. The derivations use only contextual equation symmetry, cut, and renaming already established above.

\begin{theorem}[Six-rule equivalence]\label{cr:six}
For every canonical native sequent,
\[
\Gamma\vdash_{\Knorm}A
\quad\Longleftrightarrow\quad
\Gamma\vdash_{\Cnorm}A.
\]
\end{theorem}
\begin{proof}
The preceding constructions provide every rule of \(\Cnorm\) inside \(\Knorm\). Conversely, a WI instance is exactly one of the two unrestricted introductions, and the crossed rules are ordinary transport or QE followed by derived equality symmetry when their conclusion is an equation. Hyp, \(R_+\), and DE are common. Transform finite derivations recursively. No numerical-substitution admissibility for \(\Knorm\) was presupposed; it may now be transferred through the equivalence.
\end{proof}

\paragraph{Syntactic tests under normalization.}
Normalization preserves literal-zero status of a whole term and the root constructor of a formula, so it does not itself alter an X or WI test. Arbitrary numerical substitution can alter such tests. Keeping the unrestricted equivalence and the later restricted-rule repair separate avoids confusing these two facts.

\section{Arithmetic bridge and both interpretation composites}

The full arithmetic bootstrap for \(\Rfree\) is inherited, not a new claim of this note. Appendices~\ref{nb:section} and~\ref{js:section} reproduce the pinned native finite-case repair and successor-seeded predecessor construction. Their recursor is denoted \(J\) rather than \(R\); renaming that one constructor makes it the present raw language.

Their explicit term \(P\) is proved, with only native cases, equality reasoning, and literal-equation induction, to satisfy
\[
P0=0,\qquad P(Sa)=a.
\]
No predecessor, successor cancellation, numerical equality test, ordinary selector, or compound-formula induction is assumed during that bootstrap. The native finite-case appendix is required: theorem-level numerical cases alone would not justify arbitrary contextual finite case splits before the logical macros exist.

The initialized scalar iterator is then the finite raw \(R\)-term
\[
\widehat I(n,a;u.h)
=
P\!\left(R(n,a;v.S\bigl(h[P(v)/u]\bigr))\right).
\]
Its zero computation is \(P(Sa)=a\). At a successor, \(R_+\) exposes an outer successor that \(P\) removes, leaving exactly
\[
\widehat I(Sn,a;u.h)=h[\widehat I(n,a;u.h)/u].
\]
The unrestricted selector \(C(n,a,b):=\widehat I(n,a;u.b)\) is now available. The established native scalar bridge reconstructs all intuitionistic first-order connectives, actual-equality reflection, quantifier-free induction, pairing, and every fixed primitive-recursive definition. That bridge is separately formalized in the pinned repository; its validity is inherited here.

Conversely, interpret a raw successor-seeded description by
\[
R(n,a;u.h)\longmapsto I(n,Sa;u.h).
\]
The two raw computations are instances of the scalar computations. For the two numerical interpretation composites, first identify successor: its definition \(R(0,a;u.0)\) translates to \(I(0,Sa;u.0)=Sa\). Then proceed through each fixed description in dependency order. After earlier descriptions have been identified, the description and its double translation have the same initial value and the same update. Literal-equation induction on the equation comparing their trajectories proves equality throughout the entire numerical domain. The initializer and passive parameters remain independent and fixed during this induction, and are instantiated only afterward. Body equality is used through the derived trajectory comparison, not as an uncharged recursor-congruence axiom.

The native formula round trips use the established logical macros once \(C\) is available. The \(Q\) and \(D\) readings are those already displayed; their inverse translations are interderivable by case introduction, DE, and QE. Formula congruence and the numerical identities extend the round trips to every formula.

Combining these inherited arithmetic and logical results with Theorems~\ref{cr:canonical-equiv} and~\ref{cr:six} gives the precise conclusion:
\[
\Knorm\ \text{is equiconsistent and bi-interpretable with full first-order }\ipra.
\]
Both numerical domains are whole and interpreted equality is actual equality. The numerical composite isomorphisms have graph \(x=y\). On raw syntax the composite through canonical syntax is provably, rather than literally, the identity; on canonical syntax normalization is literally idempotent.

\section{What has and has not become simpler}

There are six explicit inference patterns instead of seven, and no independent successor constructor. The omitted zero computation has become a canonical-syntax law:
\[
R(0,a;u.h)\equiv R(0,a;u.0).
\]
That law must be enforced by every substitution and comparison used by a complete checker. It has not disappeared from the trusted specification.

For the current compact encodings, a numerical recursor and a \(Q\)-formula both have arity three, so their opener must distinguish numerical depth before applying this normalization. Testing the original counter instead of the substituted one is unsound. Merely rejecting noncanonical input also fails: internally generated substitution instances must be normalized.

Accordingly this theorem establishes a mathematical change of presentation and a smaller explicit rule inventory. The separately measured matching strict implementation is also smaller; no absolute or compact-checker minimum for this new calculus follows. A fair measurement must include canonicalization, syntax and scope handling, decoding, every inference rule, and the public input/output driver. Any separate implementation improvement for the old scalar calculus must not be attributed to this new canonical \(R\) calculus.

\appendix
\section*{Provenance of the following inherited appendices}
The following two source files are reproduced from
\file{Research/successor_seed_completion/theory/}
at the pinned commit: local \file{successor_seed_prelogic.tex} is upstream
\file{prelogic.tex}, and local \file{successor_seed_bootstrap.tex} is upstream
\file{bootstrap.tex}. Their mathematics is inherited. They make explicit the arithmetic result invoked by the new normalization proof. Their own references to earlier shared-continuation work are retained.
% Native bridge audit and the pre-selector finite-case repair.
% Requires amsmath, amssymb, amsthm and theorem/lemma environments.

\section{Native bootstrap: cases before a full selector}
\label{nb:section}

This section checks a dependency that matters for the new scalar basis.
Before predecessor is available, the partial selector
\[
 B(n,a,b)=J(n,a;u.b),\qquad B(0,a,b)=Sa,\quad B(Sn,a,b)=b
\]
is not the unrestricted selector $C(0,a,b)=a$, $C(Sn,a,b)=b$.
Accordingly, one must not silently invoke the later ordinary-connective
macros during the arithmetic bootstrap.  In particular, the theorem rule
for numerical cases is not by itself a rule for eliminating equations
$t=0$ and $t=Sz$ from a retained context.  The finite-case lemma below
supplies exactly the stronger form used by the bootstrap.

We use the native Eq/$Q$/$D$ presentation with hypothesis use, full
formula transport, $Q$ elimination, paired elimination, and the two
paired introductions.  In notation,
\[
 Q_yF=\text{the clause taking a proof of }F(0)\text{ to }F(t),
\]
\[
 \frac{\Gamma,F(t,0)\vdash F(t,Sz)}
      {\Gamma\vdash D_{x,y}F(x,y)},\qquad
 \frac{\Gamma,E(t,z)\vdash E(t,Sz)}
      {\Gamma\vdash D_{x,y}E(x,y)}.
\]
The second matrix is a literal equation.  The usual freshness conditions
are retained, and paired elimination opens $Q_yF(a,y)$ at a fresh $a$.
Computation equations may be read as leaves, or obtained from their
constant generated-assumption instances and hypothesis use.

\subsection{Resources available without a numerical selector}

Choose $v$ fresh for $t$ and put $r=J(S0,0;v.t)$.  Computation gives
$r=t$.  Transport this equation along itself with motive $w=t$ to obtain
$t=t$.  The familiar equality motives then give symmetry, transitivity
and numerical congruence; this uses neither predecessor nor successor
injectivity.

For any native formula template $F$, case introduction with an unused
witness coordinate gives $D_{x,y}F(y)$ from
$\Gamma\vdash F(Sz)$, ignoring the retained $F(0)$.  Paired elimination
opens $Q_yF(y)$; a supplied proof of $F(0)$ followed by $Q$ elimination
at $t$ gives
\[
 \frac{\Gamma\vdash F(0)\qquad\Gamma\vdash F(Sz)}
      {\Gamma\vdash F(t)}. \tag{Native NC}\label{nb:nc}
\]
The identical construction with progress introduction, using its local
$E(z)$, gives contextual literal-equation induction from a base and step.
These arguments use no arithmetic selector or ordinary logical connective.

For reference, $Q_yF$ and $D_{x,y}F$, with $x$ unused, are interderivable
already: use case introduction and $Q$ elimination in the forward
direction, and paired elimination followed by hypothesis use in the
reverse direction.  Thus the additional $Q$ constructor presents the
same internal clauses as the Eq/$D$ formulation.

Even equation explosion is available at this stage.  Under $0=1$,
congruence gives $B(0,0,t)=B(1,0,t)$; computation gives $1=t$.
Consequently every equation follows.  External induction on fixed native
formula syntax, using case introduction (and the unused-witness
representation for $Q$), gives explosion for every native formula.
This is an external proof transformation, not compound-formula induction
inside the arithmetic theory.

\subsection{Functional finite ranges do not assume branch equations}

Say that a particular term $b$ has the \emph{functional range rule}
$\operatorname{FR}_R(b)$, for a finite list of numerals $R$, when the
following metarule is derived for every native formula template $K$:
\[
 \bigl(\Gamma\vdash K(r)\text{ for each }r\in R\bigr)
 \quad\Longrightarrow\quad\Gamma\vdash K(b).
 \tag{FR}\label{nb:fr}
\]
There are no added equations $b=r$ in these premises.  The retained
context and every fixed parameter of $K$ remain unchanged.  In
particular, they may mention the original term $b$.

Use the fixed terms $Z,N,j,M$ defined in
Section~\ref{js:cycles}, with $Z(t)=B(t,0,0)$. The present proof
checks their logical prerequisites; their finite computation identities
are verified directly from those definitions.

For $Z(t)$, apply Native NC to the fresh template $w\mapsto K(Z(w))$.
The zero branch is $K(1)$ and every successor branch is $K(0)$, by
computation and transport.  Hence $\operatorname{FR}_{\{0,1\}}(Z(t))$
holds, and therefore so does the same rule for $N(t)$.
For $B(t,1,0)$ the identical construction has endpoints two and zero.
For
\[
 j(t)=B\bigl(t,0,B(Z(N(t)),1,0)\bigr)
\]
use NC on the outer control: the zero branch is one, and the successor
branch has the preceding two-value rule.  This proves
$\operatorname{FR}_{\{0,1,2\}}(j(t))$, and hence the same rule for
$M(t)$, since $M$ is an outer application of $j$.

The inner parity iterate $U(t)=J(t,0;u.Z(u))$ has a two-value range
rule as well.  At counter zero it is one; at a successor it is a value
of $Z$.  Apply NC on the counter and the already proved range rule of
$Z$.  No induction with a disjunctive invariant is involved.

These constructions justify the cycle equations before contextual
finite-case elimination is used.  For example, $Z(Z(U(t)))=U(t)$ follows
by applying its two-value rule to the equation template $Z(Z(w))=w$.
The equations
\[
 M(0)=0,\quad M(1)=1,\quad M(2)=2,\qquad M(Sw)=j(M(w))
\]
and $j(0)=1$, $j(1)=2$, $j(2)=0$ then follow from the fixed definitions.
Their functional range proofs are independent of the ramp invariant.

\begin{lemma}[Successor cancellation on the three finite values]
\label{nb:finite-cancellation}
Suppose $b$ and $c$ have functional range rules whose values are contained
in $\{0,1,2\}$.  In any retained context,
\[
                         Sb=Sc\quad\Longrightarrow\quad b=c
\]
is derivable, before global successor injectivity is known.
\end{lemma}
\begin{proof}
Apply the two functional range rules to the equation templates
$M(w)=w$ and $j(j(j(w)))=w$.  Their required numeral instances are fixed
computations.  Thus $M(b)=b$, $M(c)=c$, $j^3(b)=b$ and $j^3(c)=c$
are already theorems in the retained context.
Apply $M$ to $Sb=Sc$ and use $M(Sw)=j(M(w))$ to get $j(b)=j(c)$.
Apply $j$ twice and use the two cycle identities to obtain $b=c$.
Every step is ordinary numerical congruence or transport.  In particular,
no cancellation of an arbitrary successor was inserted.
\end{proof}

\begin{lemma}[Contextual finite-case elimination for equations]
\label{nb:guarded-cases}
Let $b$ have $\operatorname{FR}_R(b)$, where $R$ is contained in
$\{0,1,2\}$, and let $F$ be the literal equation $u=v$.
If, for each $r\in R$, there is a derivation
\[
                         \Gamma,b=r\vdash u=v,
\]
then $\Gamma\vdash u=v$ is derivable.  The context may contain arbitrary
native formulas, including an induction hypothesis that depends on $b$.
\end{lemma}
\begin{proof}
Define a native formula template, with both bound variables fresh,
\[
 K(x):=D_{a,y}\bigl(B(y,x,u)=B(y,b,v)\bigr),
\]
where $a$ is unused.  Fix $r\in R$.  To introduce $K(r)$ by the case
choice with witness zero, its local zero assumption reduces by
computation to $Sr=Sb$.  Lemma~\ref{nb:finite-cancellation}, with the
constant $r$ and the known finite-range term $b$, yields $b=r$.
Cut in the given branch derivation.  Its conclusion $u=v$ is exactly
the successor matrix, up to the computations of $B$.  Thus
$\Gamma\vdash K(r)$.

Apply the functional range rule of $b$ to these closed branch proofs.
It gives $\Gamma\vdash K(b)$, whose matrix is
$B(y,b,u)=B(y,b,v)$.  Paired elimination opens its $Q$ clause.
At zero its two sides are $Sb$, so reflexivity and computation supply
the base.  Specializing the clause at one gives $u=v$.  Close paired
elimination.  All temporary coordinates are fresh and all temporary
assumptions have been discharged.
\end{proof}

The derivation is a finite guard construction, not a primitive implication
connective.  It relies only on the partial selector, computations,
transport, hypotheses, case introduction, paired elimination and $Q$
elimination.  Neither predecessor nor the unrestricted selector occurs.

\subsection{Audit of the bootstrap and completion of the bridge}

The dependency order is now explicit.  First establish functional range
rules and the two finite cycles using only Native NC and equations.
Then establish finite successor cancellation and
Lemma~\ref{nb:guarded-cases}.  The finite splits in the ramp induction,
the affine scaling induction and the factor-map calculation are instances
of that lemma, with $b=N(n)$ or another displayed finite-range term and
with a literal-equation conclusion.  Retained induction hypotheses are
therefore legitimate.  The saturation identity $K(c,Su)=S K(c,u)$ is
proved directly by NC with no branch assumptions; it does not require
an unrestricted contextual split on arbitrary $c$.  No step in this order
uses the later ordinary logical macros.

After the corrected predecessor $P$ is constructed, the initializer
compiler gives the ordinary scalar $I$.  Consequently
\[
 C(n,a,b):=I(n,a;u.b)
\]
has the unrestricted selector endpoints.  The inherited gate, padding,
formula-selector and ordinary-connective constructions may now be applied
in their stated order.  Binder congruence itself uses only case
introduction, paired elimination, transport and cut; it is already valid
before this final stage.

For clarity, the ordinary interpretations of the native binders are
\[
 (Q_yB)^\circ=B^\circ(0)\to\forall y\,B^\circ(y),\qquad
 (D_{x,y}B)^\circ=
 \exists x\bigl(B^\circ(x,0)\to\forall y\,B^\circ(x,y)\bigr).
\]
Once the logical macros are available, the native return translations
are proved by the same explicit constructions:
\[
 Q_yB\dashv\vdash B(0)\Rightarrow_*\forall_*y\,B(y),
\]
\[
 D_{x,y}B\dashv\vdash
 \exists_*x\bigl(B(x,0)\Rightarrow_*\forall_*y\,B(x,y)\bigr).
\]
Forwards, an opened $Q$ clause under its base supplies the arbitrary
instance needed by universal introduction.  Backwards, a displayed
implication and universal supply the successor matrix for case
introduction; an unused witness gives the $Q$ case.  The existential
witness is retained in the $D$ case.  External structural induction on
finite formulas, together with the endpoint proofs for the ordinary
macros, yields both formula composites.  Falsity uses the explicit
ordinary separation bridge $(0=1)\to\bot$.

Finally, the inherited scalar completion and fixed-description PR
compiler are used only after the logical interface and initialized $I$
have been recovered.  Recurrence uniqueness gives the numerical
composites on the whole original domain.  The domain is unchanged,
interpreted equality is actual equality, and the numerical isomorphism
graph is $x=y$.  Thus the finite-case repair closes the pre-selector
logical dependency; it does not replace the separately supplied complete
arithmetic and logical interpretations by standard-model calculations.

% Successor-seeded scalar iteration: a self-contained arithmetic bootstrap.
% Requires amsmath, amssymb, amsthm and theorem/lemma/proposition environments.
% No document wrapper. All induction invariants below are literal equations.

\section{A successor-seeded scalar basis with two computation schemas}
\label{js:section}

This section completes the explicitly unresolved one-variable
successor-seeded candidate in the supplied shared-continuation
report~\cite{shared}, in its section ``A genuinely smaller numerical
grammar: the precise unresolved step.''  The candidate basis and its
conditional initializer compiler are inherited from that report.
Unlike the counter-aware successor recursor, its update has just one
changing input.  The initializer remains a separate argument, but the
iterator starts at its successor.  The missing result is a numerical
predecessor term with a whole-domain proof, obtained here by adapting
the supplied scalar $2/3$ factor-map construction to shifted arithmetic
and correcting a saturated predecessor at three small arguments.

\subsection{Syntax, computations, and the precise resources used}

The terms are finite staged expressions
\[
 t ::= x\mid 0\mid J(n,a;u.h).
\]
Only $h$ binds $u$; its other free variables are passive parameters.
Define the abbreviation
\[
 Sa:=J(0,a;u.0),\qquad 1:=S0,
\]
and use standard numerals as finite strings of this abbreviation.
There are exactly two numerical computation schemas:
\begin{align}
 J(0,a;u.h)&=Sa,\tag{J0}\label{js:j0}\\
 J(Sn,a;u.h)&=h[J(n,a;u.h)/u].\tag{JS}\label{js:js}
\end{align}
The first schema, with the abbreviation expanded, says that
$J(0,a;u.h)=J(0,a;u.0)$.  It does not presuppose a separate successor
constructor.  No function variables or uniform evaluator are present.

The bootstrap uses equality replacement, numerical cases
\[
 \frac{\Gamma\vdash F(0)\qquad\Gamma\vdash F(Sz)}
      {\Gamma\vdash F(t)},
 \tag{NC}\label{js:nc}
\]
and contextual induction only on a literal equation $E$:
\[
 \frac{\Gamma\vdash E(0)\qquad\Gamma,E(z)\vdash E(Sz)}
      {\Gamma\vdash E(t)}.
 \tag{EI}\label{js:ei}
\]
The ordinary eigenvariable conditions apply.  In an ordinary
intuitionistic first-order presentation, NC follows from coverage
$\forall x\,(x=0\vee\exists y\,x=Sy)$.  In the native paired-binder
presentation it is the already available case rule.  No use of
successor injectivity, predecessor, pairing, an equality characteristic,
or induction on a compound formula is made before predecessor is
constructed below.

The finite range splits used below require the distinction proved in
Section~\ref{nb:section}. Functional range elimination follows from NC;
contextual finite cases with retained induction hypotheses use the finite
cancellation and guard construction of Lemma~\ref{nb:guarded-cases}.
The cycle identities are established first by functional range elimination
alone. Thus a displayed finite range is not an uncharged disjunction or
an assumption of the ordinary-connective macros.

We shall use the following admissible consequence of EI.  If two fixed
updates are pointwise equal, and the displayed initial values agree,
their iterates agree at every counter.  Induct on the literal equation
between the two iterates: J0 supplies its base and JS supplies its step.
This is a derived uniqueness argument, not a primitive extensionality
rule for update bodies.  All passive arguments are kept fixed while
performing the induction.

\subsection{A partial selector and shifted addition}

Put
\[
 B(n,a,b):=J(n,a;u.b),\qquad Z(n):=B(n,0,0),
 \qquad a\oplus b:=J(b,a;u.Su).
\]
Their computation endpoints are
\[
 B(0,a,b)=Sa,\quad B(Sn,a,b)=b,\qquad
 Z(0)=1,\quad Z(Sn)=0,
\]
\[
 a\oplus0=Sa,\qquad a\oplus Sb=S(a\oplus b).
\]
The operation $\oplus$ has standard meaning $a+b+1$.  Ordinary
unshifted addition has not been defined and is not used in the proof.
Literal-equation induction proves
\[
 0\oplus a=Sa,\quad Sa\oplus b=S(a\oplus b),\quad
 a\oplus b=b\oplus a,\quad
 (a\oplus b)\oplus c=a\oplus(b\oplus c).
 \tag{Shifted algebra}\label{js:oplus}
\]
For the first two equations induct on the second argument.  For
commutativity induct on $b$, using the first two equations.  For
associativity induct on $c$; at zero both sides are
$S(a\oplus b)$, and the successor step applies $S$ to the previous
equality.  These proofs do not cancel a successor.

In particular, the following are genuine available terms:
\[
 D(a):=S^2(a\oplus a),\qquad
 T(a):=S^4((a\oplus a)\oplus a),\qquad
 R(a):=S^3(a\oplus a).
 \tag{Affine terms}\label{js:affine}
\]
Their standard meanings are respectively $2a+3$, $3a+6$, and
$2a+4$.  The proof uses their displayed term definitions.  From
\eqref{js:oplus} one obtains
\[
 D(0)=3,\quad D(Sa)=S^2D(a),\qquad
 T(0)=6,\quad T(Sa)=S^3T(a).
 \tag{Affine shifts}\label{js:affine-shifts}
\]

\subsection{Two finite cycles}\label{js:cycles}

Define
\[
 N(n):=Z\bigl(J(n,0;u.Z(u))\bigr).
\]
The inner iterator is always zero or one: at zero it is one, and at a
successor it is a value of $Z$.  The endpoints of $Z$ and finite cases
therefore give
\[
 N(0)=0,\quad N(Sn)=Z(N(n)),\quad N(S^2n)=N(n),
 \qquad N(n)\in\{0,1\}.
 \tag{Parity cycle}\label{js:parity}
\]
This asserts a proved two-state recurrence, not a pre-existing parity
operation.  For example, if $K(n)=J(n,0;u.Z(u))$, then
$N(Sn)=Z(Z(K(n)))=K(n)=Z(N(n))$ by its two cases.

Next put
\[
 j(u):=B\bigl(u,0,B(Z(N(u)),1,0)\bigr),\qquad
 M(n):=j\bigl(j(J(n,0;u.j(u)))\bigr).
\]
Finite computations give
\[
 j(0)=1,\qquad j(1)=2,\qquad j(2)=0.
\]
The range of $j$ is contained in $\{0,1,2\}$: its zero branch is
one, and its successor branch is zero or two according to $N(u)$.
Consequently the definitions and these three computations give
\[
 M(0)=0,\quad M(Sn)=j(M(n)),\quad M(S^3n)=M(n),
 \qquad M(n)\in\{0,1,2\}.
 \tag{Three-cycle}\label{js:three-cycle}
\]
Here $M(0)=j(j(1))=0$, and the successor equation follows by
associativity of composing the fixed term $j$ with itself.  The period
three follows from its three possible values.  No numerical division
algorithm has been assumed.

\subsection{The shifted ramp and factor map}

Set
\[
 G(n):=J(n,0;u.u\oplus M(S^2u)).
\]
Its first values are
\[
 G(0)=1,\ G(1)=2,\ G(2)=4,\ G(3)=5,\
 G(4)=7,\ G(5)=8,\ G(6)=10.
\]

\begin{lemma}[Ramp identities]\label{js:ramp}
The preceding theory proves
\[
 M(G(n))=S(N(n)),\qquad G(S^2n)=S^3G(n),
 \tag{Ramp}\label{js:ramp-identities}
\]
and admits the following two finite-case calculations:
\[
 \begin{array}{c|cc}
 &N(n)=0&N(n)=1\\ \hline
 G(Sn)&S G(n)&S^2G(n)\\
 G(S^3n)&S^4G(n)&S^5G(n).
 \end{array}
 \tag{Ramp cases}\label{js:ramp-cases}
\]
\end{lemma}
\begin{proof}
Induct on the single equation $M(G(n))=S(N(n))$.
Its base is $M(1)=1$.  If $N(n)=0$, the induction equation gives
$M(G(n))=1$, hence $M(S^2G(n))=0$ by the three-cycle.  The next
state is $G(n)\oplus0=S G(n)$, whose $M$-value is two, as
required.  If $N(n)=1$, the current $M$-value is two and
$M(S^2G(n))=1$.  The next state is
$G(n)\oplus1=S^2G(n)$, whose $M$-value is one.  This proves
the invariant and the first line of the case table.  Applying that line
twice, with the parity cycle, gives the two-step equation; applying it
three times gives the last line.
\end{proof}

\begin{lemma}[Affine identities]\label{js:affine-lemmas}
Throughout the entire numerical domain,
\[
 N(D(a))=1,\qquad N(T(a))=N(a),\qquad
 S G(D(a))=T(a),
 \tag{Affine parity and ramp}\label{js:affine-parity}
\]
\[
 G(T(a))=T(G(a))\oplus N(a).
 \tag{Ramp scaling}\label{js:ramp-scaling}
\]
\end{lemma}
\begin{proof}
For $N(D(a))=1$, use EI, $D(0)=3$, and the two-step parity
period.  For $N(T(a))=N(a)$ use EI, $T(0)=6$, and
$T(Sa)=S^3T(a)$; three parity steps toggle both sides once.

For $S G(D(a))=T(a)$ use EI on that literal equation.  The base
is $S G(3)=6$.  The successor step uses
$D(Sa)=S^2D(a)$ and $G(S^2v)=S^3G(v)$, so it applies three
successors to the previous equation, just as $T$ does.

For \eqref{js:ramp-scaling}, the base is
$G(6)=10=T(1)\oplus0$.  Retain its equation as the induction
hypothesis, and split $N(a)$.  In the zero case, $N(T(a))=0$,
so the new left side is $S^4G(T(a))$.  The old right side is
$T(G(a))\oplus0=S T(G(a))$, while the new right side is
\[
 T(SG(a))\oplus1=S^5T(G(a)).
\]
Thus both sides have changed by four successors.  In the one case,
the new left side is $S^5G(T(a))$.  The old right side is
$T(G(a))\oplus1=S^2T(G(a))$, while the new right side is
\[
 T(S^2G(a))\oplus0=S^7T(G(a)).
\]
Thus both sides have changed by five successors.  These calculations
use only \eqref{js:oplus}, \eqref{js:affine-shifts}, and the ramp
case table; no successor has been cancelled.
\end{proof}

Define the factor map
\[
 H(u):=B(Z(N(u)),G(u),u).
\]
Its two cases are
\[
 N(u)=0\ \Longrightarrow\ H(u)=u,\qquad
 N(u)=1\ \Longrightarrow\ H(u)=S G(u).
\]

\paragraph{Why these shifts are useful.}
For intuition on standard natural numbers only, $G(u)=\lfloor3u/2\rfloor+1$.
Thus $H$ fixes even inputs and sends an odd $u$ to $3(u+1)/2$.
After shifting the state by three, an active step multiplies it by $3/2$.
The affine maps $D(a)=2a+3$ and $T(a)=3a+6$ are precisely doubling
and tripling in this shifted coordinate: $D(a)+3=2(a+3)$ and
$T(a)+3=3(a+3)$. Starting at two makes the unshifted base even and
fixed by $H$, while its shifted value is five. The later seed sequence
therefore has standard values $5\cdot2^x-3$. This motivates a comparator
that tests when its powers of two have been exhausted. The proof below
uses the displayed $J$ terms and literal equations; division, exponentiation,
and this standard-value description are not prerequisites.

\begin{lemma}[The factor-map equations]\label{js:factor}
The following global equations hold:
\[
 H(2)=2,\qquad H(D(a))=T(a),\qquad H(T(a))=T(H(a)).
 \tag{Factor}\label{js:factor-equations}
\]
\end{lemma}
\begin{proof}
The first equation follows from $N(2)=0$.  Since $N(D(a))=1$,
the second is $S G(D(a))=T(a)$.
For the last equation split $N(a)$, noting that $N(T(a))=N(a)$.
In the zero case both sides equal $T(a)$.  In the one case,
\eqref{js:ramp-scaling} gives
\[
 H(T(a))=S G(T(a))
 =S\bigl(T(G(a))\oplus1\bigr)
 =S^3T(G(a))=T(SG(a))=T(H(a)).
\]
\end{proof}

\subsection{A comparator without stripping an initializer}

Put
\[
 A(x):=J(x,0;u.R(u)),\qquad E(x):=S A(x),
\]
\[
 F(x,y):=J(y,A(x);u.H(u)),\qquad
 L(x,y):=Z(N(F(x,y))).
\]
By direct computation and shifted algebra,
\[
 E(0)=2,\qquad E(Sx)=D(E(x)),\qquad F(x,0)=E(x).
\]
The definition of $F$ is the decisive initializer alignment:
its successor-seeded initial value is exactly $E(x)=S A(x)$.
There is no hidden use of predecessor to initialize it at $E(x)$.

\begin{lemma}[Comparator equations]\label{js:comparator}
The term $L$ is zero or one and satisfies
\[
 L(0,y)=1,\qquad L(Sx,0)=0,\qquad
 L(Sx,Sy)=L(x,y).
 \tag{Comparator}\label{js:comparator-equations}
\]
\end{lemma}
\begin{proof}
Its range follows from the range of $Z$.  EI on $y$ proves
$F(0,y)=2$, since $E(0)=2$ and $H(2)=2$; hence $L(0,y)=1$.
Also $F(Sx,0)=E(Sx)=D(E(x))$ has parity one, giving
$L(Sx,0)=0$.

For the remaining equation first prove, by EI on $y$ with $x$ fixed,
\[
 F(Sx,Sy)=T(F(x,y)).
\]
At zero, the left side is
$H(E(Sx))=H(D(E(x)))=T(E(x))$.  In the successor step both
iterates apply $H$, and $H\circ T=T\circ H$ rewrites the result
into $T(F(x,Sy))$.  Applying $Z\circ N$ and using
$N(T(v))=N(v)$ gives the desired comparator equation.

This proof avoids cancellation in an equation between two successor
initializers.  Such cancellation would be unavailable at this stage.
\end{proof}

\subsection{Saturation and a finite correction give true predecessor}

Define conditional increment by
\[
 K(c,u):=B(c,u,u),\qquad
 B_m(u):=K(Z(L(Su,m)),u),\qquad
 V(t,m):=J(t,0;u.B_m(u)).
\]
Thus $K(0,u)=Su$ and $K(Sv,u)=u$.  The counter $m$ is a passive
parameter of the iteration defining $V$.

\begin{lemma}[Shifted saturation]\label{js:saturation}
The theory proves
\[
 V(St,S^2k)=S V(t,Sk),\qquad
 V(n,Sn)=Sn,\qquad
 V(Sn,Sn)=V(S^2n,Sn)=Sn.
 \tag{Saturation}\label{js:saturation-equations}
\]
\end{lemma}
\begin{proof}
Numerical cases on $c$ give $K(c,Su)=S K(c,u)$.  Consequently
the comparator equations give
\[
 B_{Sm}(Su)=S B_m(u),\qquad B_{S^2k}(1)=2.
\]
EI on $t$, with $k$ fixed, now proves the first displayed equation.
At zero its two sides equal two.  At a successor, apply the defining
update to the induction equation and then use the conjugacy
$B_{S^2k}(Su)=S B_{Sk}(u)$.

EI gives $L(Sn,n)=0$ directly from the comparator equations.
A separate EI proves $V(n,Sn)=Sn$: the base is $V(0,1)=1$,
and the successor step uses the previously proved uniform shift at
$t=k=n$.  This is substitution into a completed lemma, not a change
of a passive parameter inside an active induction hypothesis.
Finally
\[
 B_{Sn}(Sn)=K(Z(L(S^2n,Sn)),Sn)
           =K(Z(L(Sn,n)),Sn)=K(1,Sn)=Sn.
\]
One and then two more iterator steps give the last two equations.
\end{proof}

Put
\[
 Q(m):=J\bigl(m,0;u.K(Z(L(S^2u,m)),u)\bigr).
\]
This term is a capped predecessor rather than the desired predecessor:
\[
 Q(0)=Q(1)=1,\qquad Q(S^2k)=Sk.
 \tag{Capped predecessor}\label{js:capped}
\]
The first two equations are direct computations.  For the last,
\[
 L(S^2u,S^2k)=L(u,k)=L(Su,Sk)
\]
identifies its update with $B_{Sk}(u)$.  The derived iterator uniqueness
lemma, proved at an independent counter $t$ with $k$ fixed, gives
\[
 J(t,0;u.K(Z(L(S^2u,S^2k)),u))=V(t,Sk).
\]
Substitute $t=S^2k$ and apply \eqref{js:saturation-equations}.

Now define one fixed correction term
\[
 \boxed{\quad
 P(m):=B\bigl(Z(L(3,m)),\ Q(Q(m)),\ L(2,m)\bigr).
 \quad}\tag{Predecessor}\label{js:predecessor}
\]
There is no arbitrary selector hidden in this expression.  Its outer
operation is the original partial selector $B$, whose zero branch adds
exactly one successor.

\begin{theorem}[Predecessor and initializer recovery]\label{js:recovery}
The two-schema successor-seeded scalar theory proves
\[
 P(0)=0,\qquad P(Sn)=n.
\]
For every fixed update $h$, the finite term
\[
 I(n,a;v.h):=
 P\bigl(J(n,a;u.S(h[Pu/v]))\bigr)
 \tag{Initializer compiler}\label{js:initializer}
\]
satisfies
\[
 I(0,a;v.h)=a,\qquad I(Sn,a;v.h)=h[I(n,a;v.h)/v].
\]
\end{theorem}
\begin{proof}
The comparator equations give the three direct values
\[
 P(0)=0,\qquad P(1)=0,\qquad P(2)=1.
\]
Indeed, at these inputs $L(3,m)=0$, so the outer $B$ selects
$L(2,m)$, whose three values are zero, zero, one.
At $m=S^3k$, one has $L(3,m)=1$.  Thus the zero-control branch
of $B$, together with \eqref{js:capped}, gives
\[
 P(S^3k)=S Q(Q(S^3k))=S Q(S^2k)=S^2k.
\]
Numerical cases twice on $n$ now give $P(Sn)=n$ from these displayed
equations.  Alternatively, two literal-equation inductions, first on
$P(S^2n)=Sn$ and then on $P(Sn)=n$, give the same conclusion.
No earlier cancellation or injectivity was used.  Successor injectivity
can now be derived by applying $P$ to an equality of successors.

For the initializer compiler, J0 and the proved predecessor equation
give the base $P(Sa)=a$.  At a successor, JS exposes the term
$S(h[P(J(n,a;u.S(h[Pu/v])))/v])$.  The outer $P$ removes its
displayed successor and leaves exactly $h[I(n,a;v.h)/v]$.
There is no invariant about the internal state to prove: the displayed
successor in the update is sufficient.
\end{proof}

\subsection{Consequences and the bi-interpretation boundary}

The construction supplies the full initialized scalar basis $0,S,I$.
In particular the unrestricted selector is now the derived term
$C(n,a,b):=I(n,a;u.b)$, with its ordinary endpoints $a$ and $b$.
The supplied scalar arithmetic development can therefore be used from
this point onward to obtain actual-equality reflection, induction on
arbitrary quantifier-free formulas, pairing and projections, and the
compiler for every fixed general primitive-recursion description.
Those later arguments are not premises of this bootstrap.

In ordinary iPRA, each fixed $J$ description is interpreted by primitive
recursion with seed $Sa$ and its displayed scalar update.  Both
computation schemas, coverage, separation, and literal-equation induction
are then provable.  Conversely \eqref{js:initializer}, followed by the
existing scalar compiler, interprets every fixed iPRA operation.
For each fixed description, literal-equation induction proves uniqueness
of its defining recurrence after earlier subdescriptions have been
identified.  Applied to both composite translations, this identifies
each numerical operation with its double translation throughout the
entire original domain.  Both domain formulas are $x=x$ and both
numerical isomorphism graphs are $x=y$.

For the ordinary presentations, full intuitionistic logic is preserved.
For the native paired-binder presentation, its ordinary logical macros
and the formula return-translation proof must still be included; they
apply once the derived selector $C$ is available.  This establishes
bi-interpretability only when that logical bridge and the scalar
arithmetic completion accompany the present argument.  It does not
replace them by agreement merely on the standard natural numbers.

The arithmetic advance is completion of the predecessor obligation
left open for this exact candidate in the shared-continuation
report~\cite{shared}: the affine factor-map alignment, capped
predecessor, and finite correction \eqref{js:predecessor} supply its
missing proof.  The candidate grammar and initializer compiler are not
new.  Finite cycles, a $2/3$ factor map, and saturation are adapted from
the supplied ordinary scalar predecessor construction.  No global
literature priority is asserted.

The now-completed basis has two primitive numerical constructors and
two computation schemas, with a single update variable.  A smaller
complete proof checker remains a separate implementation claim, which
must charge the actual decoder, scope checks, logical rules, and driver.


\begin{thebibliography}{9}
\bibitem{shared}
ProveIt, ConciseIPRA research collection, successor-seeded and shared-continuation research lineage, as documented in
\file{Research/successor_seed_completion/README.txt},
commit \texttt{8eb71c1e570e07ef45ead686628e77f18ef34798}.
The inherited appendices identify their earlier antecedents; no priority claim is made in the present note.
\end{thebibliography}
\end{document}

